	\section{Trabajo a futuro}
	
	Este trabajo representa una primera aproximación al reconocimiento de la LSM utilizando visión por computadora en dispositivos móviles. Sin embargo, aún existen múltiples áreas de mejora y líneas de investigación que podrían incrementar la precisión, robustez, escalabilidad y aplicabilidad del prototipo desarrollado.
	
	Una de las principales líneas de trabajo a futuro consiste en realizar todo el proceso de adquisición del conjunto de datos directamente desde dispositivos Android. El conjunto de datos fue generado y procesado en computadora, mientras que la inferencia final se ejecuta en un teléfono móvil. Aunque se replicó el mismo preprocesamiento durante la etapa de inferencia, continúan existiendo diferencias entre ambos entornos, principalmente debido a variaciones en la orientación de la cámara, resolución, procesamiento interno de imágenes, distancia focal, iluminación y comportamiento específico de MediaPipe en Android.
	
	Se propone desarrollar un módulo de captura nativo dentro de la propia aplicación móvil que permita almacenar directamente las coordenadas de los landmarks detectados por MediaPipe. Estas coordenadas podrían guardarse en archivos de texto plano, JSON o CSV para posteriormente ser transferidas a una computadora y utilizadas durante el entrenamiento del modelo. Este enfoque permitiría generar un conjunto de datos más representativo de las condiciones reales de uso, reduciendo el desfase entre entrenamiento e inferencia. Posteriormente el modelo entrenado podría convertirse nuevamente a TensorFlow Lite y reintegrarse a Android para evaluar posibles mejoras en estabilidad, precisión y generalización.
	
	Otra línea importante de investigación consiste en ampliar el vocabulario de señas soportadas por el sistema. El prototipo actual se enfoca principalmente en señas estáticas correspondientes a letras del alfabeto de la LSM sin embargo, existen múltiples señas adicionales que podrían incorporarse gradualmente, incluyendo números, palabras comunes, saludos, expresiones básicas y vocabulario cotidiano. Para ello, sería necesario construir nuevos conjuntos de datos balanceados, considerando diferentes usuarios, tamaños de mano, distancias y condiciones de iluminación. Asimismo, sería recomendable analizar aquellas señas con configuraciones manuales similares, ya que representan uno de los principales retos para los modelos de clasificación.
	
	Resulta relevante continuar incrementando la diversidad del conjunto de datos. Como trabajo futuro, se propone realizar capturas utilizando distintos dispositivos móviles, diferentes cámaras frontales y traseras, múltiples resoluciones y escenarios reales de uso. También sería conveniente incorporar usuarios de distintas edades y características físicas, así como considerar variaciones relacionadas con mano dominante, velocidad de movimiento y ángulos de ejecución de las señas. Esto permitiría construir un sistema con mayor capacidad de generalización y menos dependiente de condiciones controladas.
	
	Otra posible mejora consiste en experimentar con técnicas adicionales de aumento y transformación de datos. Aunque el uso de landmarks reduce considerablemente la dependencia del fondo y la iluminación, aún podrían evaluarse estrategias como perturbaciones geométricas, ruido controlado, rotaciones leves, escalamiento dinámico o interpolación temporal de secuencias. Estas técnicas podrían ayudar a mejorar la robustez del modelo frente a variaciones naturales durante la captura.
	
	Es importante continuar evaluando distintos modelos de aprendizaje automático y aprendizaje profundo. Durante el desarrollo del presente trabajo se analizaron múltiples enfoques, incluyendo CNNs y modelos basados en landmarks; sin embargo, aún podrían explorarse arquitecturas híbridas o modelos más compactos específicamente optimizados para dispositivos móviles. También podría analizarse el uso de técnicas de cuantización, poda de redes neuronales o aceleración mediante NNAPI y GPU para reducir la latencia y el consumo computacional.
	
	Un experimento que resulta relevante consiste en el reconocimiento de señas dinámicas. El prototipo actual se centra únicamente en señas estáticas; no obstante, muchas señas dentro de la LSM dependen del movimiento y de la trayectoria temporal de la mano. Para abordar este problema, podría analizarse el uso de modelos secuenciales como redes LSTM, GRU o arquitecturas Transformer capaces de procesar secuencias de landmarks en lugar de una única postura estática.
	
	La implementación de este tipo de modelos requeriría definir una metodología específica para la captura de secuencias, incluyendo la cantidad de frames por muestra, la duración temporal de cada seña, criterios de inicio y finalización, así como mecanismos de sincronización y segmentación temporal. Adicionalmente, sería necesario evaluar cuidadosamente el impacto de estos modelos sobre el rendimiento en dispositivos móviles, particularmente en términos de memoria, consumo energético y tiempo de inferencia.
	
	En términos de experiencia de usuario, otra línea futura consiste en mejorar la interfaz gráfica y los mecanismos de interacción. Podrían incorporarse ayudas visuales más avanzadas para el posicionamiento correcto de la mano, retroalimentación háptica o auditiva, indicadores de estabilidad de predicción y módulos interactivos para aprendizaje de señas. Asimismo, podrían desarrollarse pruebas formales de usabilidad y accesibilidad con usuarios reales para evaluar aspectos relacionados con eficiencia, satisfacción y facilidad de uso.
	
	Se recomienda continuar realizando pruebas experimentales en distintos escenarios reales de operación, evaluando métricas relacionadas con precisión, capacidad de generalización, latencia de inferencia, estabilidad de reconocimiento y consumo de recursos. Esto permitiría identificar las principales limitaciones del sistema y establecer una base sólida para futuras versiones más robustas y escalables del prototipo.
	\newpage
	
